\section{Verification methodology}\label{sec:verification}

Three activities are kept apart (\cref{sec:steps}): solid cross-code
comparison (\cref{sec:ver:solid}), implementation checks
(\cref{sec:ver:impl}), and FSI solution verification
(\cref{sec:ver:error,sec:ver:mesh,sec:ver:time,sec:ver:robust,sec:ver:cons}).

\subsection{Solid cross-code comparison}\label{sec:ver:solid}
Problem~S (AR25 Steps~0A, 0B, 1) with the original AR25 boundary conditions,
on M1--M4, compared component-wise with the nine AR25 solutions at
$\vect p_0$, $\vect p_1$, and with $\RED$ against the AR25 ensemble mean. The
AR25 solutions agree to within about \SI{0.5}{mm}, a fraction of their
element size (3--\SI{5}{mm}); this spread, and the resolution of each
participant, bound what agreement can mean. The existing regression test of
PR~\#571 (Step~1, M1, to \SI{0.05}{s}, checking $p_{0,x}\in[0.3,0.6]$\,mm)
is retained and extended to peak contraction and relaxation. If the cap
decision of \cref{sec:geometry} constrains the solid base, a solid-only
control with the same constraint is added and reported separately; the
modified problem is not called AR25's.

\subsection{Implementation checks}\label{sec:ver:impl}
\begin{itemize}
\item constitutive response on prescribed homogeneous deformation states
  (uniaxial along $\f$, $\s$, $\vect n_0$; simple shear in the $fs$ plane;
  pure dilatation), with and without $\tau$ and $\dot\EE$, against an
  independent evaluation of \cref{eq:S,eq:Psi};
\item orthonormality of $(\f,\s,\vect n_0)$ and its behaviour at the apex;
\item independent integration of \cref{eq:tau,eq:p} and comparison with the
  released tables;
\item ALE geometric conservation on a compatible moving-mesh test (a
  uniform flow on a domain with a prescribed, divergence-compatible
  boundary motion; a uniform flow is not compatible with the deforming
  no-slip LV walls, so this is a separate test, not a property of the
  production case);
\item conservation of force and work across the interface mapping;
\item a simpler established coupled problem (e.g.\ the pressure-driven
  elastic tube) with the same coupling settings \TODO{choose}.
\end{itemize}

\subsection{Error estimation}\label{sec:ver:error}
Meshes are indexed from fine to coarse, $h_1<h_2<h_3$, so that M4, M3, M2
are levels 1, 2, 3 of a triplet. For a scalar quantity $\phi$ on three
levels with refinement ratios $r_{21}=h_2/h_1$, $r_{32}=h_3/h_2$, the
procedure of \citet{Celik2008} gives the observed order $p_{\mathrm{obs}}$
(solved iteratively when $r_{21}\ne r_{32}$), the extrapolated value
$\phi_{\mathrm{ext}}$ and the fine-grid convergence index
\begin{equation}
\mathrm{GCI}_{21} = \frac{1.25\,|\phi_1-\phi_2|/|\phi_1|}{r_{21}^{\,p_{\mathrm{obs}}}-1}.
\end{equation}
Each triplet is first classified as monotonically convergent, oscillatory or
divergent from the sign of $(\phi_3-\phi_2)/(\phi_2-\phi_1)$; the GCI is
reported only for monotonic convergence, and $p_{\mathrm{obs}}$, the ratios
and the classification are reported in every case. Near-zero quantities
use an absolute error with a fixed scale. For time series the relative
error against the finest solution,
$e_{L^2} = \|\phi_h-\phi_{h_1}\|_{L^2(0,T)}/\|\phi_{h_1}\|_{L^2(0,T)}$, is used;
it is distinct from $\RED$. GCI is an estimated discretisation uncertainty
under the asymptotic-range assumption, not a bound on the error.

\paragraph{Scalar QoIs}
\begin{itemize}
\item Problem F: $V$ at the time of peak pressure and at $t=1$\,s; peak
  $\Delta p_{\mathrm{IVPD}}$ and its time; peak $E_k$ and its time;
  $u_x(\vect p_0)$ and $u_x(\vect p_1)$ at $t=0.45$\,s; peak $\Gamma$;
\item Problem C: net ejected volume at $t=1$\,s and gross outward volume;
  peak outflow rate and its time; peak $\Delta p_{\mathrm{IVPD}}$; peak
  $E_k$; $u_x(\vect p_0)$ and $u_x(\vect p_1)$ at $t=0.45$\,s;
\item both: $\max|J-1|$, $\max|R_V|/V_0$.
\end{itemize}

\subsection{Spatial convergence}\label{sec:ver:mesh}
At a fixed small $\dt$ and tight iterative tolerances (so that temporal and
iterative errors are shown to be smaller than the differences measured):
(i) paired refinement of solid and fluid, M1--M4; (ii) solid refined with
the fluid held at its finest affordable level, and (iii) the converse, to
attribute the error to each domain. Solid and fluid $h$ are reported
separately and never collapsed into one. The interface mapping, the
curvature representation and the fibre field change with refinement; their
effect is noted.

\subsection{Temporal convergence}\label{sec:ver:time}
On the finest affordable mesh,
$\dt \in \{\num{1e-3},\num{5e-4},\num{2.5e-4},\num{1.25e-4},\num{6.25e-5}\}$\,s
\TODO{adjust from pilot cost}, with $\epsilon_{\mathrm{FSI}}$ and the inner
tolerances tightened so that they do not contaminate the temporal error;
maximum Courant number reported. The observed temporal order is measured
and compared with the formal orders of the components (BDF2 for the solid
inertia, first order for the viscous term, Euler or BDF2 for the fluid);
a second-order component scheme does not imply second-order coupled
results when start-up, mesh fluxes, coupling lag, interpolation or loose
tolerances dominate. This study measures the temporal convergence of one
spatially discrete problem; it does not establish resolution independence.

\subsection{Sensitivity to numerical settings and robustness}\label{sec:ver:robust}
Starting from the baseline (\cref{tab:settings}) and a tightly converged
reference (finest mesh and $\dt$, tightest tolerances), one setting is
varied at a time on M2, and the most influential on M3. A run is a success
only if it completes, every time step meets the coupling criterion, every
inner solve converges, $|R_V|/V_0$ and the mesh-quality limits are
satisfied. For each run: success, deviation of the QoIs from the converged
reference (not from the baseline), FSI iterations per step and wall-clock
time.
\begin{enumerate}
\item \textbf{Solid incompressibility}: mixed $\uu$--$p$ vs.\ pure
  displacement with penalty $\kappa$; $\max|J-1|$.
\item \textbf{Solid momentum stabilisation}: $c_{\uu}\in\{1,10,100,1000\}$
  \TODO{range}; the trade-off between checkerboarding and artificial
  stiffening, and its dependence on $h$ (the term should vanish as
  $h\to0$).
\item \textbf{Solid pressure stabilisation}: $c_p\in\{0.1,1,10\}$, with
  $c_p^{\mathrm J}$ reported separately if it affects convergence only.
\item \textbf{Fluid time scheme}: Euler vs.\ backward.
\item \textbf{Fluid convection}: upwind, linearUpwind, linear.
\item \textbf{Mesh motion}: diffusivity; minimum cell quality at peak
  deformation.
\item \textbf{Interface condition}: Dirichlet--Neumann with and without
  the Robin-type pressure condition; residual normalisation floors over two
  orders of magnitude (expected to change the admissible residual, not the
  converged solution).
\item \textbf{Acceleration}: fixed $\omega\in\{0.3,0.5,1\}$, Aitken, IQN-ILS
  (with and without reuse of past time steps).
\item \textbf{Coupling tolerance}: $\epsilon_{\mathrm{FSI}}\in\{10^{-3},10^{-4},10^{-5},10^{-6}\}$;
  the loosest value whose QoIs agree with the tightest within the
  discretisation uncertainty is recommended.
\end{enumerate}
One-factor-at-a-time screening does not reveal interactions; a small
targeted interaction study is added for stabilisation~$\times$~mesh,
acceleration~$\times$~$\dt$, and normalisation floor~$\times$~tolerance. The
outcome is a robustness map (\cref{tab:robust}) and a recommended
configuration.

\subsection{Conservation and consistency checks}\label{sec:ver:cons}
\begin{itemize}
\item mass: $R_V(t)$ of \cref{sec:qoi} for F-p, C-p and F-Q (impermeable
  moving walls assumed; any interface leakage
  $\int_{\Gfsi}(\vv-\ww)\cdot\nn\,\mathrm da$ reported);
\item fluid energy, with every boundary $\partial\Of$ included:
  \begin{equation}
  \frac{\mathrm dE_k}{\mathrm dt} + \Phi_{\mathrm f} =
  \int_{\partial\Of}\vv\cdot(\sig_{\mathrm f}\nn)\,\mathrm da
  - \int_{\partial\Of}\tfrac12\rho_{\mathrm f}|\vv|^2(\vv-\ww)\cdot\nn\,\mathrm da,
  \end{equation}
  whose discrete residual is reported as a numerical budget error; the
  coupled budget adds solid kinetic and strain energy, viscous and dashpot
  dissipation, spring energy and active-stress power with the implemented
  sign convention;
\item interface consistency: velocity and traction jumps across $\Gfsi$ at
  convergence;
\item volume consistency: $V$ from the fluid mesh vs.\ from the closed
  solid-side surface.
\end{itemize}

\subsection{Staged run plan}\label{sec:ver:plan}
\begin{enumerate}
\item Specification gate: geometry, cap, ports, loads, pressure units,
  software provenance; machine-readable parameter file released.
\item Problem~S on M1--M4 with the original AR25 conditions; modified-base
  control if needed.
\item Implementation checks of \cref{sec:ver:impl}.
\item Pilot F-p and C-p on M2: reversals, $R_V$, mesh validity, pressure
  scales, solver convergence, cost per step.
\item Iterative convergence: tighten fluid, solid and interface criteria
  until the QoIs stop changing relative to the targeted discretisation
  differences.
\item Resolution studies (\cref{sec:ver:mesh,sec:ver:time}).
\item Targeted sensitivity (\cref{sec:ver:robust}).
\item Extensions: F-Q once its volume compatibility and convergence are
  shown; C-R only if fully specified.
\end{enumerate}
At $\dt=\SI{5e-5}{s}$ a one-second run is \num{20000} steps; with 30
coupling iterations per step this is $6\times10^{5}$ coupled iterations
before inner solves, so the campaign is sized from the pilot cost
measurements, with hardware, MPI ranks, memory, output overhead and timing
variability reported.
